\documentclass[xcolor={dvipsnames,svgnames},14pt]{beamer}
\usetheme[progressbar=none]{Singapore}
\setbeamercolor*{structure}{fg=cyan}
\setbeamercolor*{frametitle}{fg=black}
\setbeamercolor*{title}{fg=black}
\usepackage{hyperref}
\usepackage{multicol}
\hypersetup{
    colorlinks=true,
    linkcolor=blue,
    filecolor=magenta,      
    urlcolor=blue,
    }

\makeatletter
\let\beamer@writeslidentry@miniframeson=\beamer@writeslidentry
\def\beamer@writeslidentry@miniframesoff{%
  \expandafter\beamer@ifempty\expandafter{\beamer@framestartpage}{}% does not happen normally
  {%else
    % removed \addtocontents commands
    \clearpage\beamer@notesactions%
  }
}
\newcommand*{\miniframeson}{\let\beamer@writeslidentry=\beamer@writeslidentry@miniframeson}
\newcommand*{\miniframesoff}{\let\beamer@writeslidentry=\beamer@writeslidentry@miniframesoff}
\makeatother

\miniframesoff

\setbeamertemplate{navigation symbols}{%
    \usebeamerfont{footline}%
    \usebeamercolor[fg]{footline}%
    \hspace{1em}%
    \insertframenumber/\inserttotalframenumber
}
\setbeamercolor{footline}{fg=black}
\setbeamerfont{footline}{series=\bfseries}

\usepackage{amsmath, amsthm, amssymb}
\usepackage{graphicx, verbatim}
\usepackage{multirow}
\usepackage{xcolor}
\DeclareRobustCommand{\stirling}{\genfrac\{\}{0pt}{}}


%Information to be included in the title page:
\title{Two-Sample Mean Inference}
\date{}
\subtitle{Estimating a difference in means}

\begin{document}

\frame{\titlepage}

\begin{frame}
Download the section 14 .Rmd handout to \\ 
\texttt{STAT240/lecture/sect14-two-means}. \\~\\

Material in this section is covered by Chapter 13 on
the notes website.
\\~\\

Download the file \texttt{cat\_breeds\_clean.txt} to \texttt{STAT240/data}.
\end{frame}

\section{Equal variances T}

\begin{frame}
Continue working with the cat sleep data. \\~\\

What is the change in the average sleep time for fixed male ragdoll cats versus intact male ragdoll cats? \\~\\
\end{frame}

\begin{frame}
Inference for $\mu_1 - \mu_2$: \\~\\ \begin{itemize}
\item Point estimate: $\bar{X}_1 - \bar{X}_2$
\item Standard error:
$$\sqrt{\frac{S^2_1}{n_1} + \frac{S^2_2}{n_2}}$$ 
\\~\\
\item Model: T distribution \\~\\
\end{itemize}
\end{frame}

\begin{frame}
Assume the groups have two different populations: \begin{align*}
X_{1,i} & \sim D_1(\mu_1,\; \sigma_1) \\
X_{2,i} & \sim D_2(\mu_2,\; \sigma_2)
\end{align*}
and populations 1 and 2 are independent. \\~\\
We allow the distributions to differ completely.  All that matters is $\bar{X}_1$ and $\bar{X}_2$ are approximately normal.
\end{frame}

\begin{frame}
$$\text{point estimate } \pm \text{ critical value } \times \text{ standard error}$$
\begin{itemize}
\item For $\mu_1 - \mu_2$, the point estimate is $\bar{x}_1 - \bar{x}_2$
\item The standard error combines $\sigma_1$ and $\sigma_2$
\end{itemize}
\end{frame}

\begin{frame}
The standard error is
$$\sqrt{se_1^2 + se_2^2} \;=\; \sqrt{\frac{\sigma^2_1}{n_1} + \frac{\sigma^2_2}{n_2}}$$
This gives sampling distribution
$$\bar{X}_1 - \bar{X}_2 \quad\dot{\sim} \quad N\Big(\mu_1 - \mu_2,\;\; \sqrt{\frac{\sigma^2_1}{n_1} + \frac{\sigma^2_2}{n_2}}\Big)$$
\end{frame}


\begin{frame}
This suggests a normal critical value.  However, we need to estimate the se.
$$\hat{se}(\bar{X}_1 - \bar{X}_2) \;=\; \sqrt{\frac{s^2_1}{n_1} + \frac{s^2_2}{n_2}}$$
\\~\\
Again, we use a T distribution for our CI and hypothesis test.
\end{frame}


\begin{frame}
There are two different types of two-sample T test:  \begin{itemize}
\item Equal variances T (often more powerful)
\item Unequal variances T (more general) \\~\\
\end{itemize}
We can assume equal variances when the sample SDs are within a factor of 2. 
$$0.5 < \frac{s_1}{s_2} < 2$$
\end{frame}

\begin{frame}
When we estimate $\sigma_1 = \sigma_2 = \sigma$, we can simplify the standard error.

$$se \;=\; \sigma \sqrt{\frac{1}{n_1} + \frac{1}{n_2}}$$ \\~\\

Estimate $\sigma$ with the \textbf{pooled sd} $s_p$.
\end{frame}

\begin{frame}
The pooled sd is a weighted average of the groups.  
$$s_p \;=\; \sqrt{\frac{(n_{1}-1)s_{1}^2 \;+\; (n_{2}-1)s_{2}^2}{n_{1} + n_{2} - 2}}$$ 
The weights are $n_1-1$ and $n_2-1$. \\~\\
In the cats example, the pooled sd is about 2.477.
\end{frame}

\begin{frame}
$$s_p \;=\; \sqrt{\frac{(n_{1}-1)s_{1}^2 \;+\; (n_{2}-1)s_{2}^2}{n_{1} + n_{2} - 2}}$$ 
\begin{itemize}
\item Is $s_p$ more similar to $s_1$ or $s_2$?
\item When is $s_p$ an exact average of $s_1$ and $s_2$?
\item[$\bigstar$] Why do we square the terms, then take an average, then square root?  Why not just average $s_1$ and $s_2$?
\end{itemize}
\end{frame}

\begin{frame}
Now, we can estimate the standard error.
$$\hat{se} \;=\; s_p \sqrt{\frac{1}{n_1} + \frac{1}{n_2}}$$ \\~\\
Lastly, we use a T with $n_1 + n_2 - 2$ degrees of freedom for the critical value.
\end{frame}

\begin{frame}
We are 90\% confident that the difference in sleep time is within (-0.91, 0.35).  More generally, we have

$$\bar{X}_1 - \bar{X}_2 \;\pm\; t_{\alpha/2, n_1+n_2-2}\cdot  s_p\sqrt{\frac{1}{n_1} + \frac{1}{n_2}}$$
\end{frame}

\begin{frame}
We can also perform a hypothesis test for the difference in means. Let's do a 10\% level test of the difference in sleep times: 
$$H_0: \mu_1 - \mu_2 = 0 \quad \text{versus} \quad H_A: \mu_1 - \mu_2 \neq 0$$
\end{frame}

\begin{frame}
The test statistic (assuming equal variances) is 
$$T \;=\; \frac{\bar{X}_1 - \bar{X}_2 - 0}{s_p \sqrt{\frac{1}{n_1} + \frac{1}{n_2}}}$$
and follows a $T$ with $n_1 + n_2 - 2$ degrees of freedom if $H_0$ is true.
\end{frame}

\begin{frame}
In general, we can test for any mean difference $\mu_0$:
$$H_0: \mu_1 - \mu_2 = \mu_0 \quad \text{versus} \quad H_A: \mu_1 - \mu_2 \neq \mu_0$$
with test statistic
$$T \;=\; \frac{\bar{X}_1 - \bar{X}_2 - \mu_0}{s_p \sqrt{\frac{1}{n_1} + \frac{1}{n_2}}}$$
\end{frame}

\section{Unequal variances T}

\begin{frame}
What if we aren't willing to make the equal variances assumption?  We can still perform a similar T analysis using the more general standard error:
$$\hat{se}(\bar{X}_1 - \bar{X}_2) \;=\; \sqrt{\frac{s^2_1}{n_1} + \frac{s^2_2}{n_2}}$$
Let's compare intact male cats and kittens.
\end{frame}

\begin{frame}
The T is approximate, and the df is messy.
$$w \;=\; \frac{\Big(\frac{s^2_1}{n_1} + \frac{s^2_2}{n_2}\Big)^2}{\frac{(s^2_1/n_1)^2}{n_1-1} + \frac{(s^2_2/n_2)^2}{n_2-1}}$$
\\~\\
This is the \textbf{Welch} procedure.  The df are ``penalized" according to $s_1$ and $s_2$.
\end{frame}

\begin{frame}
The formula for the CI is
$$\bar{X}_1 - \bar{X}_2 \;\pm\; t_{\alpha/2, w}\cdot \sqrt{\frac{s^2_1}{n_1} + \frac{s^2_2}{n_2}}$$
\\~\\
We are 95\% confident that the true difference is within (-1.416, 1.961).  \texttt{t.test()} confirms this.
\end{frame}


\begin{frame}
Are adult cats more sleepy than kittens?
$$H_0: \mu_1 - \mu_2 = 0 \quad \text{versus}\quad H_A: \mu_1 - \mu_2 > 0$$
\\~\\
Let's complete this test with $\alpha = 0.05$.
\end{frame}

\begin{frame}
Our observed test statistic is 
$$T \;=\; \frac{\bar{X}_1 - \bar{X}_2 - 0}{\sqrt{\frac{s^2_1}{n_1} + \frac{s^2_2}{n_2}}}$$ 
\\~\\
If $H_0$ is true, $T$ has a T distribution with $w$ degrees of freedom. \\~\\ \begin{itemize}
\item Complete the test of the one-sided hypotheses.
\end{itemize}
\end{frame}

\begin{frame}
\begin{center}
\includegraphics[scale=0.7]{t_null.png}
\end{center}
\end{frame}

\begin{frame}
For hypotheses
$$H_0: \mu_1 - \mu_2 = 0 \quad \text{versus}\quad H_A: \mu_1 - \mu_2 > 0$$
we have p-value
$$P(T_w > 0.316) = 0.997$$
and fail to reject $H_0$.
\end{frame}

\begin{frame}
\begin{center}
\includegraphics[scale=0.7]{t_pval.png}
\end{center}
\end{frame}

\begin{frame}
$\alpha = 0.05$ rejection region:
\vspace{-5pt}
\begin{center}
\includegraphics[scale=0.65]{t_rej_region.png}
\end{center}
\end{frame}

\begin{frame}
P-value calculation: \begin{itemize}
\item If the alternative is $H_A: \text{diff} < \mu_0$, the p-value is the area \textit{below} the test statistic. 
\item If the alternative is $H_A: \text{diff} > \mu_0$, the p-value is the area \textit{above} the test statistic. 
\item If the alternative is $H_A: \text{diff} \neq \mu_0$, the p-value is $2\times$ the area \textit{outside} of the test statistic.
\end{itemize}
\end{frame}

\begin{frame}
What if we assumed equal variances?   In general: \begin{itemize}
\item If the variances are truly equal, the Welch T test is a bit less powerful, but is still accurate. 
\item If the variances are truly different, the equal variances test can make incorrect conclusions.  \\~\\
\end{itemize}
It is safer to assume unequal variances.
\end{frame}

\section{Paired T test}

\begin{frame}
Sometimes, two samples of data are \textbf{paired}. \\~\\

Do children grow taller from age 13 to 14?  It doesn't make sense to take independent samples of 13 and 14 year olds. \\~\\

Instead, measure the same individuals at 13 and 14.
\end{frame}

\begin{frame}
Five individuals had their heights measured at age 13 and 14.
\begin{center}
\begin{tabular}{|l||c|c|c|c|c|}
\hline
Height at 13 & 44.1 & 59.0 & 65.9 & 58.7 & 49.3 \\
\hline
Height at 14 & 46.3 & 60.5 & 68.2 & 59.4 & 50.6 \\
\hline
\end{tabular}
\end{center}
\vspace{5pt}
The two samples are \textit{not independent}, so we can't use one of the previous methods.
\end{frame}

\begin{frame}
Instead, study the \textit{differences}.
\begin{center}
\begin{tabular}{|l||c|c|c|c|c|}
\hline
Height at 13 & 44.1 & 59.0 & 65.9 & 58.7 & 49.3 \\
\hline
Height at 14 & 46.3 & 60.5 & 68.2 & 59.4 & 50.6 \\
\hline
\hline
Difference & 2.2 & 1.5 & 2.3 & 0.7 & 1.3 \\
\hline
\end{tabular}
\end{center}
\vspace{5pt}
We are back in a one-sample setting!  Parameter:
$$\mu_{diff} \quad \text{instead of}\quad \mu_1 - \mu_2$$
\end{frame}

\begin{frame}
Assume the differences are independent draws from a population:
$$diff_i \;\sim\; D(\mu, \sigma)$$
We then have hypotheses
$$H_0: \mu_{diff} = 0 \quad \text{versus} \quad H_A: \mu_{diff} > 0$$
\end{frame}

\begin{frame}
We use the usual one-sample T test statistic
$$T \;=\; \frac{\bar{diff} - \mu}{S/\sqrt{n}}$$ \\~\\
which follows a T with $n - 1$ df if $H_0$ is true. \\~\\

Let's complete the test with $\alpha = 0.05$.
\end{frame}

\begin{frame}
We have a large positive test statistic and a small p-value, so we reject $H_0$. \\~\\

This is confirmed in \texttt{t.test()}. \\~\\

If we had done a two independent sample test, we get the opposite result!
\end{frame}


\end{document}





